% ==== IV ====
\yoschapter{IV}{The Hahn-Banach Theorems}
\label{chap:IV}

In a Hilbert space, we can introduce the notion of orthogonal coordinates through an orthogonal base, and these coordinates are the values of bounded linear functionals defined by the vectors of the base. This suggests that we consider continuous linear functionals, in a linear topological space, as generalized coordinates of the space. To ensure the existence of non-trivial continuous linear functionals in a general locally convex linear topological space, we must rely upon the Hahn-Banach extension theorems. 


\yossection{1}{The Hahn-Banach Extension Theorem in Real Linear Spaces}
\label{sec:IV.1}

\begin{theoremu}[\textsc{Hahn} {\cite{hahn90}}, \textsc{Banach} {\cite{banach9}}] Let $X$ be a real linear space and let $p(x)$ be a real-valued function defined on $X$ satisfying the conditions:
\begin{align}
p (x + y)& \leqq p (x) + p (y) \quad (\text{ subadditivity }), \label{eq:IV.1.1} \\
p (\alpha x)& = \alpha p (x) \quad \text{ for } \quad \alpha \geqq 0. \label{eq:IV.1.2}
\end{align}

Let $M$ be a real linear subspace of $X$ and $f_0$ a real-valued linear functional defined on $M$: 
\begin{equation}\label{eq:IV.1.3}
f _{0} (\alpha x + \beta y) = \alpha f _{0} (x) + \beta f _{0} (y) \quad \text{ for } \quad x, y \in M \text{ and real } \alpha , \beta .
\end{equation}

Let $f_0$ satisfy $f_0(x) \leqq p(x)$ on $M$. Then there exists a real-valued linear functional $F$ defined on $X$ such that \begin{enumerate}[i)]
\item $F$ is an extension of $f_0$, i.e., $F(x) = f_0(x)$ for all $x \in M$
\item $F(x) \leqq p(x)$ on $X$.
\end{enumerate}


\end{theoremu}
\begin{proof}
 Suppose first that $X$ is spanned by $M$ and an element $x_0 \bar{\in} M$, that is, suppose that 
\[
X = \{x = m + \alpha x _{0}; m \in M, \alpha \text{ real } \}. \]

Since $x_0 \bar{\in} M$, the above representation of $x \in X$ in the form $x = m + \alpha x_0$ is unique. It follows that, if, for any real number $c$, we set 
\[
F (x) = F (m + \alpha x _{0}) = f _{0} (m) + \alpha c,
\]
then $F$ is a real linear functional on $X$ which is an extension of $f_0$. We have to choose $c$ such that $F(x) \leqq p(x)$, that is, $f_0(m) + \alpha c \leqq p(m + \alpha x_0)$. This condition is equivalent to the following two conditions: 
\[
\begin{gathered}
f _{0} (m / \alpha) + c \leqq p (x _{0} + m / \alpha) \quad \text{ for } \quad \alpha > 0,
\\
f _{0} (m / (- \alpha)) - c \leqq p (- x _{0} + m / (- \alpha)) \quad \text{ for } \quad \alpha < 0.
\end{gathered}
\]

To satisfy these conditions; we shall choose $c$ such that 
\[
f _{0} \left(m ^{\prime}\right) - p \left(m ^{\prime} - x _{0}\right) \leqq c \leqq p \left(m ^{\prime \prime} + x _{0}\right) - f _{0} \left(m ^{\prime \prime}\right) \text{ for all } m ^{\prime}, m ^{\prime \prime} \in M.
\]

Such a choice of $c$ is possible since 
\[
\begin{aligned} f _{0} (m ^{\prime}) + f _{0} (m ^{\prime \prime}) & = f _{0} (m ^{\prime} + m ^{\prime \prime}) \leqq p (m ^{\prime} + m ^{\prime \prime}) = p (m ^{\prime} - x _{0} + m ^{\prime \prime} + x _{0}) \\ & \leqq p (m ^{\prime} - x _{0}) + p (m ^{\prime \prime} + x _{0}): \end{aligned}
\]
we have only to choose $c$ between the two numbers 
\[
\sup\limits _ {m ^{\prime} \in M} [ f _{0} (m ^{\prime}) - p (m ^{\prime} - x _{0}) ] \quad \text{ and } \quad \inf\limits _ {m ^{\prime \prime} \in M} [ p (m ^{\prime \prime} + x _{0}) - f _{0} (m ^{\prime \prime}) ].
\]

Consider now the family of all real linear extensions $g$ of $f_0$ for which the inequality $g(x) \leqq p(x)$ holds for all $x$ in the domain of $g$. We make this family into a partially ordered family by defining $h > g$ to mean that $h$ is an extension of $g$. Then Zorn's Lemma ensures the existence of a maximal linear extension $g$ of $f_0$ for which the inequality $g(x) \leqq p(x)$ holds for all $x$ in the domain of $g$. We have to prove that the domain $D(g)$ of $g$ coincides with $X$ itself. If it does not, we obtain, taking $D(g)$ as $M$ and $g$ as $f_0$, a proper extension $F$ of $g$ which satisfies $F(x) \leqq p(x)$ for all $x$ in the domain of $F$, contrary to the maximality of the linear extension $g$. 
\end{proof}

\begin{corollaryu} Given a functional $p(x)$ defined on a real linear space $X$ such that \eqref{eq:IV.1.1} and \eqref{eq:IV.1.2} are satisfied. Then there exists a linear functional $f$ defined on $X$ such that 
\begin{equation}\label{eq:IV.1.4}
- p (- x) \leqq f (x) \leqq p (x).
\end{equation}
\end{corollaryu}

\begin{proof}
 Take any point $x_0 \in X$ and define $M = \{x; x = \alpha x_0, \alpha \text{ real }\}$. Set $f_0(\alpha x_0) = \alpha p(x_0)$. Then $f_0$ is a real linear functional defined on $M$. We have $f_0(x) \leqq p(x)$ on $M$. In fact, $\alpha p(x_0) = p(\alpha x_0)$ if $\alpha > 0$, and if $\alpha < 0$, we have $\alpha p(x_0) \leqq -\alpha p(-x_0) = p(\alpha x_0)$ by $0 = p(0) \leqq p(x_0) + p(-x_0)$. Thus there exists a linear functional $f$ defined on $X$ such that $f(x) = f_0(x)$ on $M$ and $f(x) \leqq p(x)$ on $X$. Since $-f(x) = f(-x) \leqq p(-x)$, we obtain $-p(-x) \leqq f(x) \leqq p(x)$. 
\end{proof}


\yossection{2}{The Generalized Limit}
\label{sec:IV.2}

The notion of a sequence $\{x_{n}\}$ of a countable number of elements $x_{n}$ is generalized to the notion of a directed set of elements depending on a parameter which runs through an uncountable set. The notion of the limit of a sequence of elements may be extended to the notion of the generalized limit of a directed set of elements. 

\begin{definitionu} $A$ partially ordered set $A$ of elements $\alpha, \beta, \ldots$ is called a directed set if it satisfies the condition: 

\begin{equation}\label{eq:IV.2.1}
\begin{minipage}[b]{0.92\linewidth}
For any pair $\alpha, \beta$ of elements of $A$, there exists $a\gamma \in A$ such that $\alpha < \gamma$, $\beta < \gamma$.
\end{minipage}
\end{equation}
\end{definitionu}


Let, to each point $\alpha$ of a directed set $A$, there be associated a certain set of real numbers $f(\alpha)$. Thus $f(\alpha)$ is a, not necessarily one-valued, real function defined on the directed set $A$. We write 
\[
\lim\limits _{\alpha \in A} f (\alpha) = a \quad (a \text{ is a real number })
\]
if, for any $\varepsilon > 0$, there exists an $\alpha_0 \in A$ such that $\alpha_0 < \alpha$ implies $|f(\alpha) - a| \leqq \varepsilon$ for all possible values of $f$ at $\alpha$. We say, in such a case, that the value $a$ is the generalized limit or Moore-Smith limit of $f(\alpha)$ through the directed set $A$. 

\begin{examplen} Consider a partition $\Delta$ of the real interval [$0$, $1$]:
\[
0 = t _{0} < t _{1} < \dots < t _{n} = 1.
\]

The totality $P$ of the partition of $[0,1]$ is a directed set $A$ by defining the partial order $\Delta \prec \Delta'$ as follows: If the partition $\Delta'$ is given by $0 = t_0' < t_1' < \dots < t_m' = 1$, then $\Delta \prec \Delta'$ means that $n \leqq m$ and that every $t_i$ is equal to some $t_j'$. Let $x(t)$ be a real-valued continuous function defined on $[0,1]$, and let $f(\Delta)$ be the totality of real numbers of the form 
\[
\sum_{j=0}^{n-1}(t_{j+1}-t_j)x(t_j'),
\]
where $t_j'$ is any point of $[t_j,t_{j+1}]$.

Thus $f(\Delta)$ is the totality of the Riemann sum of the function $x(t)$ pertaining to the partition $\Delta$. The value of the Riemann integral $\int_{0}^{1} x(t) \dd t$ is nothing but the generalized limit of $f(\Delta)$ through $P$. 

As to the existence of a generalized limit, we have the 

\end{examplen}
\begin{theoremn}[S. \textsc{Banach}] Let $x(\alpha)$ be a real-valued bounded function defined on a directed set $A$. The totality of such functions constitutes a real linear space $X$ by
\[
(x + y) (\alpha) = x (\alpha) + y (\alpha), (\beta x) (\alpha) = \beta x (\alpha).
\]

We can then define a linear functional, defined on $X$ and which we shall denote by $\LIM x(\alpha)$, satisfying the condition 
\[
\varliminf\limits_{\alpha \in A} x (\alpha) \leqq \operatorname * {L I M} _{\alpha \in A} x (\alpha) \leqq \varlimsup\limits_{\alpha \in A} x (\alpha),
\]
where 
\[
\varliminf\limits_{\alpha \in A} x (\alpha) = \sup\limits _{\alpha} \inf\limits _{\alpha < \beta} x (\beta), \varlimsup\limits_{\alpha \in A} x (\beta) = \inf\limits _{\alpha} \sup\limits _{\alpha < \beta} x (\beta).
\]

Therefore $\lim\limits_{\alpha \in A}x(\alpha) = \lim\limits_{\alpha \in A}x(\alpha)$ if the latter generalized limit exists. 

\end{theoremn}
\begin{proof}
 We put $p(x) = \varlimsup\limits_{\alpha \in A} x(\alpha)$. It is easy to see that this $p(x)$ satisfies the condition of the Hahn-Banach extension theorem. Hence, there exists a linear functional $f$ defined on $X$ such that $-p(-x) \leqq f(x) \leqq p(x)$ on $X$. We can easily prove that $\varliminf\limits_{\alpha \in A} x(\alpha) = -p(-x)$ so that we obtain the Theorem, by putting $\LIM_{\alpha \in A} x(\alpha) = f(x)$. 
\end{proof}


\yossection{3}{Locally Convex, Complete Linear Topological Spaces}
\label{sec:IV.3}

\begin{definitionu} As in numerical case, we may define a directed set $\{x_{\alpha}\}$ in a linear topological space $X$. $\{x_{\alpha}\}$ is said to converge to an element $x$ of 
\end{definitionu}

$X$, if, for every neighbourhood $U(x)$ of $x$, there exists an index $\alpha_{0}$ such that $x_{\alpha} \in U(x)$ for all indices $\alpha > \alpha_{0}$. A directed set $\{x_{\alpha}\}$ of $X$ is said to be fundamental, if every neighbourhood $U(0)$ of the zero vector $0$ of $X$ is assigned an index $\alpha_{0}$ such that $(x_{\alpha} - x_{\beta}) \in U(0)$ for all indices $\alpha, \beta > \alpha_{0}$. A linear topological space $X$ is said to be complete if every directed fundamental set of $X$ converges to some element $x \in X$ in the sense above. 

\begin{remarku} We can weaken the condition of the completeness, and require only that every sequence of $X$ which is fundamental as a directed set converges to an element $x \in X$; a space $X$ satisfying this condition is said to be sequentially complete. For normed linear spaces, the two definitions of completeness are equivalent. However, in the general case, not every sequentially complete space is complete. 
\end{remarku}

\begin{examplen} Let a sequence $\{f_h(x)\}$ of $\mathfrak{D}(\Omega)$ satisfy the condition $\lim\limits_{h,k\to \infty}(f_h - f_k) = 0$ in $\mathfrak{D}(\Omega)$. That is, by the Corollary of Proposition~\ref{prop:I.1.Proposition.7} in Chapter~\ref{chap:I}, \ref{sec:I.1}, we assume that there exists a compact subset $K$ of $\Omega$ for which $\supp(f_h)\subseteq K$ ( $h = 1,2,\ldots$) and $\lim\limits_{h,k\to \infty}(D^s f_h(x) - D^s f_k(x)) = 0$ uniformly on $K$ for any differential operator $D^{s}$. Then it is easy to see, by applying the Ascoli-Arzelà theorem, that there exists a function $f\in \mathfrak{D}(\Omega)$ for which $\lim\limits_{h\to \infty}D^{s}f_{h}(x) = D^{s}f(x)$ uniformly on $K$ for any differential operator $D^{s}$. Hence $\lim\limits_{h\to \infty}f_{h} = f$ in $\mathfrak{D}(\Omega)$ and so $\mathfrak{D}(\Omega)$ is sequentially complete. Similarly, we can prove that $\mathfrak{E}(\Omega)$ is also sequentially complete. 
\end{examplen}

As in the case of a normed linear space, we can prove the 

\begin{theoremu} Every locally convex linear topological space $X$ can be embedded in a locally convex, complete linear topological space, in which $X$ forms a dense subset. 
\end{theoremu}

We omit the proof. The reader is referred to the literature listed in J. A. \textsc{Dieudonné} \cite{dieudonne39}. Cf. also G. \textsc{Köthe} \cite{konishi164}. 


\yossection{4}{The Hahn-Banach Extension Theorem in Complex Linear Spaces}
\label{sec:IV.4}

\begin{theoremn}[\textsc{Bohnenblust}-\textsc{Sobczyk}] Let $X$ be a complex linear space and $p$ a semi-norm defined on $X$. Let $M$ be a complex linear subspace of $X$ and $f$ a complex linear functional defined on $M$ such that $|f(x)| \leqq p(x)$ on $M$. Then there exists a complex linear functional $F$ defined on $X$ such that \begin{enumerate}[i)]
\item $F$ is an extension of $f$
\item $|F(x)| \leqq p(x)$ on $X$.
\end{enumerate}


\end{theoremn}
\begin{proof}
 We observe that a complex linear space is also a real linear space if the scalar multiplication is restricted to real numbers. If $f(x) = g(x) + ih(x)$, where $g(x)$ and $h(x)$ are the real and imaginary parts of $f(x)$ respectively, then $g$ and $h$ are real linear functionals defined on $M$. 

Thus 
\[
\left| g (x) \right| \leqq \left| f (x) \right| \leqq p (x) \quad \text{ and } \quad \left| h (x) \right| \leqq \left| f (x) \right| \leqq p (x) \text{ on } M.
\]

Since, for each $x \in M$, 
\[
g (i x) + i h (i x) = f (i x) = i f (x) = i (g (x) + i h (x)) = - h (x) + i g (x),
\]
we have 
\[
h (x) = - g (i x) \quad \text{ for all } \quad x \in M.
\]

By the Theorem in Chapter~\ref{chap:IV}, \ref{sec:IV.1}, we can extend $g$ to a real linear functional $G$ defined on $X$ with the property that $G(x) \leqq p(x)$ on $X$. Hence $-G(x) = G(-x) \leqq p(-x) = p(x)$, and so $|G(x)| \leqq p(x)$. We define 
\[
F (x) = G (x) - i G (i x).
\]

Then, by $F(ix) = G(ix) - iG(-x) = G(ix) + iG(x) = iF(x)$, we easily see that $F$ is a complex linear functional defined on $X$. $F$ is an extension of $f$. For, $x \in M$ implies that 
\[
F (x) = G (x) - i G (i x) = g (x) - i g (i x) = g (x) + i h (x) = f (x).
\]

To prove $|F(x)| \leqq p(x)$, we write $F(x) = r e^{-i\theta}$ so that $|F(x)| = e^{i\theta} F(x) = F(e^{i\theta} x)$ is real positive; consequently $|F(x)| = |G(e^{i\theta} x)| \leqq p(e^{i\theta} x) = |e^{i\theta}| p(x) = p(x)$. 
\end{proof}


\yossection{5}{The Hahn-Banach Extension Theorem in Normed Linear Spaces}
\label{sec:IV.5}

\begin{theorem}\label{thm:IV.5.Theorem.1} Let $X$ be a normed linear space, $M$ a linear subspace of $X$ and $f_{1}$ a continuous linear functional defined on $M$. Then there exists a continuous linear functional $f$ defined on $X$ such that \begin{enumerate}[i)]
\item $f$ is an extension of $f_{1}$
\item $\|f_{1}\| = \|f\|$.
\end{enumerate}
\end{theorem}

\begin{proof}
 Set $p(x) = \|f_1\| \cdot \|x\|$. Then $p$ is a continuous semi-norm defined on $X$ such that $|f_1(x)| \leqq p(x)$ on $M$. There exists, by the Theorem in the preceding section $4$, a linear extension $f$ of $f_1$ defined on the whole space $X$ and such that $|f(x)| \leqq p(x)$. Thus $\|f\| \leqq \sup\limits_{\|x\| \leqq 1} p(x) = \|f_1\|$. On the other hand, since $f$ is an extension of $f_1$, we must have $\|f\| \geqq \|f_1\|$ and so we obtain $\|f_1\| = \|f\|$. 
\end{proof}

\yossub{An Application to Moment Problems}

\begin{theorem}\label{thm:IV.5.Theorem.2} Let $X$ be a normed linear space. Given a sequence of elements $\{x_{n}\} \subseteq X$, a sequence of complex numbers $\{\alpha_{n}\}$ and a positive number $\gamma$. Then a necessary and sufficient condition for the existence of a continuous linear functional $f$ on $X$ such that $f(x_{i}) = \alpha_{i} (i = 1, 2, \ldots)$ and $\|f\| \leqq \gamma$ is that the inequalities 
\[
\left| \sum\limits_{i=1} ^{n} \beta_{i} x _{i} \right| \leqq \gamma \left\| \sum\limits_{i=1} ^{n} \beta_{i} x _{i} \right\|
\]
hold for any choice of positive numbers $n$ and complex numbers $\beta_1, \beta_2, \ldots, \beta_n$. 
\end{theorem}

\begin{proof}
 The necessity of this condition is clear from the definition of $\|f\|$. We shall prove the sufficiency. Consider the set 
\[
X _{1} = \left\{z; z = \sum\limits_{i=1} ^{n} \beta_{i} x _{i} \quad \text{ where } n \text{ and } \beta \text{ are arbitrary } \right\}. \]

For two representations $z = \sum\limits_{i=1}^{n} \beta_i x_i = \sum\limits_{i=1}^{m} \beta_{i'} x_{i'}$ of the same element $z \in X_1$, we have, by the condition of the Theorem, 
\[
\left| \sum\limits_{i=1} ^{n} \beta_{i} \alpha_{i} - \sum\limits_{i=1} ^{m} \beta_ {i ^{\prime}} \alpha_ {i ^{\prime}} \right| \leqq \gamma \left\| \sum\limits_{i=1} ^{n} \beta_{i} x _{i} - \sum\limits_{i=1} ^{m} \beta_ {i ^{\prime}} x _ {i ^{\prime}} \right\| = 0.
\]

Thus a continuous linear functional $f_{1}$ is defined on $X_{1}$ by $f_{1}\left(\sum\limits_{i=1}^{n}\beta_{i}x_{i}\right)=\sum\limits_{i=1}^{n}\beta_{i}\alpha_{i}$. We have only to extend, by the preceding Theorem~\ref{thm:IV.5.Theorem.1}, $f_{1}$ to a continuous linear functional $f$ on $X$ with $\|f\|=\|f_{1}\|$. 
\end{proof}

\begin{remarku} As will be shown in section $9$ of this chapter, any continuous linear functional $f$ on $C[0,1]$ is representable as 
\[
f (x) = \int_{0} ^{1} x (t) m (\dd t)
\]
\end{remarku}

with a uniquely determined Baire measure $m$ on the interval [$0$, $1$]. Thus, if we take $x_{j}(t) = t^{j-1} (j = 1, 2, \ldots)$, Theorem~\ref{thm:IV.5.Theorem.2} gives the solvability condition of the moment problem: 
\[
\int_{0} ^{1} t ^{j-1} m (\dd t) = \alpha_{j} (j = 1, 2, \dots).
\]


\yossection{6}{The Existence of Non-trivial Continuous Linear Functionals}
\label{sec:IV.6}

\begin{theorem}\label{thm:IV.6.Theorem.1} Let $X$ be a real or complex linear topological space, $x_{0}$ a point of $X$ and $p(x)$ a continuous semi-norm on $X$. Then there exists a continuous linear functional $F$ on $X$ such that $F(x_{0}) = p(x_{0})$ and $|F(x)| \leqq p(x)$ on $X$. 
\end{theorem}

\begin{proof}
 Let $M$ be the set of all elements $\alpha x_{0}$, and define $f$ on $M$ by $f(\alpha x_{0}) = \alpha p(x_{0})$. Then $f$ is linear on $M$ and $|f(\alpha x_{0})| = |\alpha p(x_{0})| = p(\alpha x_{0})$ there. Thus there exists, by the Theorem in Chapter~\ref{chap:IV}, \ref{sec:IV.4}, an extension $F$ of $f$ such that $|F(x)| \leqq p(x)$ on $X$. Hence $F(x)$ is continuous at $x = 0$ with $p(x)$, and so, by the linearity of $F$, $F(x)$ is continuous at every point of $X$. 
\end{proof}

\begin{corollary}\label{cor:IV.6.Corollary.1} Let $X$ be a locally convex space and $x_{0} \neq 0$ be an element of $X$. Then there exists a continuous semi-norm $p$ on $X$ such that $p(x_{0}) \neq 0$. Thus, by Theorem~\ref{thm:IV.6.Theorem.1}, there exists a continuous linear functional $f_{0}$ on $X$ such that 
\[
f _{0} (x _{0}) = p (x _{0}) \neq 0 \quad \text{ and } \quad | f _{0} (x) | \leqq p (x) \text{ on } X.
\]
\end{corollary}

\begin{corollary}\label{cor:IV.6.Corollary.2} Let $X$ be a normed linear space and $x_{0} \neq 0$ be any element of $X$. Then there exists a continuous linear functional $f_{0}$ on $X$ such that 
\[
f _{0} (x _{0}) = \| x _{0} \| \quad \text{ and } \quad \| f _{0} \| = 1.
\]
\end{corollary}

\begin{proof}
 We take $\|x\|$ for $p(x)$ in Corollary~\ref{cor:IV.6.Corollary.1}. Thus $\|f_0\| \leqq 1$ from $|f_0(x)| \leqq \|x\|$. But, by $f_0(x_0) = \|x_0\|$, we must have the equality $\|f_0\| = 1$. 
\end{proof}

\begin{remarku} As above, we prove the following theorem by the Theorem in Chapter~\ref{chap:IV}, \ref{sec:IV.1}. 
\end{remarku}

\begin{theoremn}[Theorem~\ref{thm:IV.6.Theorem.1}']\yosnamedtag{1'}\label{thm:IV.6.Theorem.1prime} Let $X$ be a real linear topological space, $x_0$ a point of $X$ and $p(x)$ a real continuous functional on $X$ such that 
\[
p (x + y) \leqq p (x) + p (y) \quad \text{ and } \quad p (\alpha x) = \alpha p (x) \quad \text{ for } \quad \alpha \geqq 0.
\]

Then there exists a continuous real linear functional $F$ on $X$ such that $F(x_0) = p(x_0)$ and $-p(-x) \leqq F(x) \leqq p(x)$ on $X$.
\end{theoremn}

\begin{theorem}\label{thm:IV.6.Theorem.2} Let $X$ be a locally convex linear topological space. Let $M$ be a linear subspace of $X$, and $f$ a continuous linear functional on $M$. Then there exists a continuous linear functional $F$ on $X$ which is an extension of $f$. 
\end{theorem}

\begin{proof}
 Since $f$ is continuous on $M$ and $X$ is locally convex, there exists an open, convex, balanced neighbourhood of $0$, say $U$, of $X$ such that $x \in M \cap U$ implies $|f(x)| \leqq 1$. Let $p$ be the Minkowski functional of $U$. Then $p$ is a continuous semi-norm on $X$ and $U = \{x; p(x) < 1\}$. For any $x \in M$ choose $\alpha > 0$ so that $\alpha > p(x)$. Then $p(x/\alpha) < 1$ and so $|f(x/\alpha)| \leqq 1$, that is $|f(x)| \leqq \alpha$. We thus see, by letting $\alpha \downarrow p(x)$, that $|f(x)| \leqq p(x)$ on $M$. Hence, by the Theorem in Chapter~\ref{chap:IV}, \ref{sec:IV.4}, we obtain a continuous linear functional $F$ on $X$ such that $F$ is an extension of $f$ and $|F(x)| \leqq p(x)$ on $X$. 
\end{proof}

\begin{theorem}[S. \textsc{Mazur}]\label{thm:IV.6.Theorem.3} Let $X$ be a real or complex, locally convex linear topological space, and $M$ a closed convex balanced subset of $X$. Then, for any $x_0 \bar{\in} M$, there exists a continuous linear functional $f$ on $X$ such that $f_0(x_0) > 1$ and $|f_0(x)| \leqq 1$ on $M$. 
\end{theorem}

\begin{proof}
 Since $M$ is closed, there exists a convex, balanced neighbourhood $V$ of $0$ such that $M \cap (x_0 + V) = \emptyset$. Since $V$ is balanced and convex, we have $\left(M + \frac{V}{2}\right) \cap \left(x_0 + \frac{V}{2}\right) = \emptyset$. The set $\left(x_0 + \frac{V}{2}\right)$ being a neighbourhood of $x_0$, the closure $U$ of $\left(M + \frac{V}{2}\right)$ does not contain $x_0$. Since $M \ni \complement$, the closed convex balanced set $U$ is a neighbourhood of $0$, because $U$ contains $\frac{V}{2}$ as a subset. Let $p$ be the Minkowski functional of $U$. Since $U$ is closed, we have, for any $x_0 \bar{\in} U$, $p(x_0) > 1$ and $p(x) \leqq 1$, if $x \in U$. 

Thus there exists, by Corollary~\ref{cor:IV.6.Corollary.1} of Theorem~\ref{thm:IV.6.Theorem.1}, a continuous linear functional $f_0$ on $X$ such that $f_0(x_0) = p(x_0) > 1$ and $|f_0(x)| \leqq p(x)$ on $X$. Hence, in particular, $|f_0(x)| \leqq 1$ on $M$. 
\end{proof}

\begin{corollaryu} Let $M$ be a closed linear subspace of a locally convex linear topological space $X$. Then, for any $x_0 \in X - M$, there exists a continuous linear functional $f_0$ on $X$ such that $f_0(x_0) > 1$ and $f_0(x) = 0$ on $M$. Moreover, if $X$ is a normed linear space and if $\text{ dis } (x_0, M) > d$, then we may take $\|f_0\| \leqq 1 / d$. 
\end{corollaryu}

\begin{proof}
 The first part is clear from the linearity of $M$. The second part is proved by taking $U=\{x;\dist(x,M)\leqq d\}$ in the proof of Theorem~\ref{thm:IV.6.Theorem.3}. 
\end{proof}

\begin{remarku} As above, we prove the following theorem by virtue of Theorem~\ref{thm:IV.6.Theorem.1}'. 
\end{remarku}

\begin{theoremn}[Theorem~\ref{thm:IV.6.Theorem.3}' (S. \textsc{Mazur})]\yosnamedtag{3'}\label{thm:IV.6.Theorem.3prime} Let $X$ be a locally convex real linear topological space, and $M$ a closed convex subset of $X$ such that $M \ni 0$. Then, for any $x_0 \bar{\in} M$, there exists a continuous real linear functional $f_0$ on $X$ such that $f_0(x_0) > 1$ and $f_0(x) \leqq 1$ on $M$.
\end{theoremn}

\begin{theorem}[S. \textsc{Mazur}]\label{thm:IV.6.Theorem.4} Let $X$ be a locally convex linear topological space, and $M$ a convex balanced neighbourhood of $0$ of $X$. Then, for any $x_{0} \bar{\in} M$, there exists a continuous linear functional $f_{0}$ on $X$ such that 
\[
f _{0} (x _{0}) \geqq \sup\limits _{x \in M} | f _{0} (x) |.
\]
\end{theorem}

\begin{proof}
 Let $p$ be the Minkowski functional of $M$. Then $p(x_0) \geqq 1$ and $p(x) \leqq 1$ on $M$. $p$ is continuous since $M$ is a neighbourhood of $0$ of $X$. Thus there exists, by Corollary~\ref{cor:IV.6.Corollary.1} of Theorem~\ref{thm:IV.6.Theorem.1}, a continuous linear functional $f_0$ on $X$ such that 
\[
f _{0} (x _{0}) = p (x _{0}) \geqq 1 \quad \text{ and } \quad | f _{0} (x) | \leqq p (x) \leqq 1 \text{ on } M. \]
\end{proof}

\begin{theorem}[E. \textsc{Helly}]\label{thm:IV.6.Theorem.5} Let $X$ be a $B$-space, and $f_{1}, f_{2}, \ldots, f_{n}$ be a finite system of bounded linear functionals on $X$. Given $n$ numbers $\alpha_{1}, \alpha_{2}, \ldots, \alpha_{n}$. Then a necessary and sufficient condition that there exists, for each $\varepsilon > 0$, an element $x_{\varepsilon} \in X$ such that 
\[
f _{i} (x _{\varepsilon}) = \alpha_{i} \quad (i = 1, 2, \dots , n) \quad \text{ and } \quad \| x _{\varepsilon} \| \leqq \gamma + \varepsilon
\]
is that the inequality 
\[
\left| \sum\limits_{i=1} ^{n} \beta_{i} \alpha_{i} \right| \leqq \gamma \left\| \sum\limits_{i=1} ^{n} \beta_{i} f _{i} \right\|
\]
holds for any choice of $n$ numbers $\beta_1, \beta_2, \ldots, \beta_n$. 
\end{theorem}

\begin{proof}
 The necessity is clear from the definition of the norm of a continuous linear functional. We shall prove the sufficiency. We may assume, without losing the generality, that $f$'s are linearly independent; otherwise, we can work with a linearly independent subsystem of $\{f_i\}$ which spans the same subspace as the original system $\{f_i\}$. 


We consider the mapping $x \to \varphi(x) = (f_1(x), f_2(x), \ldots, f_n(x))$ of $X$ onto the Hilbert space $l^2(n)$ consisting of all vectors $x = (\xi_1, \xi_2, \ldots, \xi_n)$ normed by $\|x\| = \left(\sum\limits_{j=1}^{n} |\xi_j|^2\right)^{1/2}$. By the open mapping theorem in Chapter~\ref{chap:II}, \ref{sec:II.5}, the image $\varphi(S_\varepsilon)$ of the sphere $S_\varepsilon = \{x \in X; \|x\| \leqq \gamma + \varepsilon\}$ contains the vector $0$ of $l^2(n)$ as an interior point for every $\varepsilon > 0$. Let us suppose that $(\alpha_1, \alpha_2, \ldots, \alpha_n)$ does not belong to $\varphi(S_\varepsilon)$. Then, by Mazur's theorem given above, there exists a continuous linear functional $F$ on $l^2(n)$ such that 
\[
F \left(\left(\alpha_{1}, \alpha_{2}, \dots , \alpha_{n}\right)\right) \geqq \sup\limits _{\| x \| \leqq \gamma + \varepsilon} | F (\varphi (x)) |.
\]

Since $l^2(n)$ is a Hilbert space, the functional $F$ is given by an element $(\beta_1, \beta_2, \ldots, \beta_n) \in l^2(n)$ in such a way that $F((\alpha_1, \alpha_2, \ldots, \alpha_n)) = \sum\limits_{j=1}^{n} \alpha_j \beta_j$. Thus 
\[
\sum\limits_{j=1} ^{n} \alpha_{j} \beta_{j} \geqq \left| \sum\limits_{j=1} ^{n} f _{j} (x) \beta_{j} \right| \quad \text{ for } \quad \| x \| \leqq \gamma + \varepsilon .
\]

But the supremum of the right hand side for $\| x \| \leqq \gamma + \varepsilon$ is $= (\gamma + \varepsilon) \times \left\| \sum\limits_{j=1}^{n} f_j \beta_j \right\|$, and this contradicts the hypothesis of the Theorem.
\end{proof}

\yossection{7}{Topologies of Linear Maps}
\label{sec:IV.7}

Let $X$, $Y$ be locally convex linear topological spaces on the same scalar field (real or complex number field). We denote by $L(X, Y)$ the totality of continuous linear operators on $X$ into $Y$. $L(X, Y)$ is a linear space by 
\begin{align}
(\alpha T + \beta S) x& = \alpha T x + \beta S x \notag \\
T, S& \in L (X, Y) \notag \\
x& \in X \notag
\end{align}

We shall introduce various topologies for this linear space $L(X, Y)$. 

\begin{enumerate}[i)]
\item Simple Convergence Topology. This is the topology of convergence at each point of $X$ and thus it is defined by the family of semi-norms of the form 
\[
p (T) = p (T; x _{1}, x _{2}, \dots , x _{n}; q) = \sup\limits _{1 \leqq j \leqq n} q (T x _{j}),
\]
where $x_{1}, x_{2}, \ldots, x_{n}$ are an arbitrary finite system of elements of $X$ and $q$ an arbitrary continuous semi-norm on $Y$. $L(X, Y)$ endowed with this topology will be denoted by $L_{s}(X, Y)$. It is clearly a locally convex linear topological space.
\item Bounded Convergence Topology. This is the topology of uniform convergence on bounded sets of $X$. Thus it is defined by the family of 
semi-norms of the form 
\[
\phi (T) = \phi (T; B; q) = \sup\limits _{x \in B} q (T x)
\]
where $B$ is an arbitrary bounded set of $X$ and $q$ an arbitrary continuous semi-norm on $Y$. $L(X, Y)$ endowed with this topology will be denoted by $L_{b}(X, Y)$. It is clearly a locally convex linear topological space.
\end{enumerate}
Since any finite set of $X$ is bounded, the simple convergence topology is weaker than the bounded convergence topology, i.e., open sets of $L_{s}(X, Y)$ are also open sets of $L_{b}(X, Y)$, but not conversely. 

\begin{definition}\label{def:IV.7.Definition.1} If $X$, $Y$ are normed linear spaces, then the topology of $L_{s}(X, Y)$ is usually called the strong topology (of operators); the one of $L_{b}(X, Y)$ is called the uniform topology (of operators). 
\end{definition}

\yossub{Dual Spaces. Weak and Weak* Topologies}

\begin{definitionn}[Definition~\ref{def:IV.7.Definition.1}']\yosnamedtag{1'}\label{def:IV.7.Definition.1prime} In the special case when $Y$ is the real or complex number field topologized in the usual way, $L(X, Y)$ is called the dual space of $X$ and will be denoted by $X'$. Thus $X'$ is the set of all continuous linear functionals on $X$. The simple convergence topology is then called the weak* topology of $X'$; provided with this topology, $X'$ will sometimes be denoted by $X'_{w*}$ and we call it the weak* dual of $X$. The bounded convergence topology for $X'$ is called the strong topology of $X'$; provided with this topology, $X'$ is sometimes denoted by $X'_{s}$ and we call it the strong dual of $X$.
\end{definitionn}

\begin{definition}\label{def:IV.7.Definition.2} For any $x \in X$ and $x' \in X'$, we shall denote by $\langle x, x' \rangle$ or $x'(x)$ the value of the functional $x'$ at the point $x$. Thus the weak* topology of $X'$, i.e., the topology of $X_{w*}'$ is defined by the family of semi-norms of the form 
\[
p \left(x ^{\prime}\right) = p \left(x ^{\prime}; x _{1}, x _{2}, \dots , x _{n}\right) = \sup\limits _{1 \leqq j \leqq n} | \langle x _{j}, x ^{\prime} \rangle |,
\]
where $x_{1}, x_{2}, \ldots, x_{n}$ are an arbitrary finite system of elements of $X$. The strong topology of $X'$, i.e., the topology of $X_{s}'$ is defined by the family of semi-norms of the form 
\[
p \left(x ^{\prime}\right) = p \left(x ^{\prime}; B\right) = \sup\limits _{x \in B} | \langle x, x ^{\prime} \rangle |
\]
where $B$ is an arbitrary bounded set of $X$. 
\end{definition}

\begin{theorem}\label{thm:IV.7.Theorem.1} If $X$ is a normed linear space, then the strong dual space $X_s'$ is a $B$-space with the norm 
\[
\left\| f \right\| = \sup\limits _{\left\| x \right\| \leqq 1} | f (x) |.
\]
\end{theorem}

\begin{proof}
 Let $B$ be any bounded set of $X$. Then $\sup\limits_{b \in B} \|b\| = \beta < \infty$, and hence $\|f\| \leqq \alpha$ implies $p(f; B) = \sup\limits_{b \in B} |f(b)| \leqq \sup\limits_{\|x\| \leqq \beta} |f(x)| \leqq \alpha \beta$. 

On the other hand, the unit sphere $S = \{x; \|x\| \leqq 1\}$ of $X$ is a bounded set, and so $\|f\| = p(f; S)$. This proves that the topology of $X_s'$ is equivalent to the topology defined by the norm $\|f\|$. 

The completeness of $X_s'$ is proved as follows. Let a sequence $\{f_n\}$ of $X_s'$ satisfy $\lim\limits_{n,m \to \infty} \|f_n - f_m\| = 0$. Then, for any $x \in X$, $|f_n(x) - f_m(x)| \leqq \|f_n - f_m\| \cdot \|x\| \to 0$ (as $n, m \to \infty$), and hence a finite $\lim\limits_{n \to \infty} f_n(x) = f(x)$ exists. The linearity of $f$ is clear. The continuity of $f$ is proved by observing that $\lim\limits_{n \to \infty} f_n(x) = f(x)$ uniformly on the unit sphere $S$. Incidentally we have proved that $\lim\limits_{n \to \infty} \|f_n - f\| = 0$. 
\end{proof}

Similarly we can prove 

\begin{theorem}\label{thm:IV.7.Theorem.2} If $X$, $Y$ are normed linear spaces, then the uniform topology (of operators) $L_{b}(X, Y)$ is defined by the operator norm 
\[
\left\| T \right\| = \sup\limits _{\left\| x \right\| \leqq 1} \left\| T x \right\|.
\]
\end{theorem}

\begin{definition}\label{def:IV.7.Definition.3} We define the weak topology of a locally convex linear topological space $X$ by the family of semi-norms of the form 
\[
p (x) = p \left(x; x _{1} ^{\prime}, x _{2} ^{\prime}, \dots , x _{n} ^{\prime}\right) = \sup\limits _{1 \leqq j \leqq n} | \langle x, x _{j} ^{\prime} \rangle |,
\]
where $x_{1}^{\prime}, x_{2}^{\prime}, \ldots, x_{n}^{\prime}$ are an arbitrary finite system of elements of $X^{\prime}$. Endowed with this topology, $X$ is sometimes denoted by $X_{w}$. 
\end{definition}


\yossection{8}{The Embedding of \texorpdfstring{$X$}{X} in its Bidual Space \texorpdfstring{$X''$}{X''}}
\label{sec:IV.8}

We first prove 

\begin{theorem}[S. \textsc{Banach}]\label{thm:IV.8.Theorem.1} Let $X$ be a locally convex linear topological space, and $X'$ its dual space. A linear functional $f(x')$ on $X'$ is of the form 
\[
f (x ^{\prime}) = \langle x _{0}, x ^{\prime} \rangle
\]
with a certain $x_0 \in X$ iff $f(x')$ is continuous in the weak* topology of $X'$. 
\end{theorem}

\begin{proof}
 The ``only if'' part is clear since $\left|\langle x_{0}, x^{\prime}\rangle\right|$ is one of the seminorms defining the weak* topology of $X^{\prime}$. The ``if'' part will be proved as follows. The continuity of $f(x^{\prime})$ in the weak* topology of $X^{\prime}$ implies that there exists a finite system of points $x_{1}, x_{2}, \ldots, x_{n}$ such that $\left|f(x^{\prime})\right| \leqq \sup\limits_{1 \leqq i \leqq n} \left|\langle x_{j}, x^{\prime}\rangle\right|$. Thus 
\[
f _{i} (x ^{\prime}) = \langle x _{i}, x ^{\prime} \rangle = 0 (i = 1, 2, \dots , n) \text{ implies } f (x ^{\prime}) = 0. \]

Consider the linear map $L: X' \to l^2(n)$, defined by 
\[
L \left(x ^{\prime}\right) = \left(f _{1} (x ^{\prime}), f _{2} (x ^{\prime}), \dots , f _{n} (x ^{\prime})\right).
\]
$L(x_1') = L(x_2')$ implies $L(x_1' - x_2') = 0$ so that $f_i(x_1' - x_2') = 0 (i = 1, 2, \ldots, n)$ and hence $f(x_1' - x_2') = 0$. Hence we may define a continuous linear map $F$ defined on the linear subspace $L(X')$ of $l^2(n)$ by 
\[
F \left(L \left(x ^{\prime}\right)\right) = F \left(f _{1} (x ^{\prime}), \dots , f _{n} (x ^{\prime})\right) = f (x ^{\prime}).
\]

This map can be extended to a continuous linear functional defined on the whole space $l^{2}(n)$; the extension is possible since $l^{2}(n)$ is of finite dimension (easier than using the Hahn-Banach extension theorem in infinite dimensional linear spaces). We denote this extension by the same letter $F$. Writing 
\[
(y _{1}, y _{2}, \dots , y _{n}) = \sum\limits_{j=1} ^{n} y _{j} e _{j}, \quad \text{ where } \quad e _{i} = (0, 0, \dots , 0, 1, 0, 0, \dots , 0)
\]
with $1$ at the $j$-th coordinate, 
we easily see that 
\[
F \left(y _{1}, y _{2}, \dots , y _{n}\right) = \sum\limits_{j=1} ^{n} y _{j} \alpha_{j}, \alpha_{j} = F \left(e _{j}\right).
\]

Therefore 
\[
f \left(x ^{\prime}\right) = \sum\limits_{j=1} ^{n} \alpha_{j} f _{j} \left(x ^{\prime}\right) = \sum\limits_{j=1} ^{n} \alpha_{j} \left\langle x _{j}, x ^{\prime} \right\rangle = \left\langle \sum\limits_{j=1} ^{n} \alpha_{j} x _{j}, x ^{\prime} \right\rangle .
\]
\end{proof}
\begin{corollaryu} Each $x_0 \in X$ defines a continuous linear functional $f_0(x')$ on $X_s'$ by $f_0(x') = \langle x_0, x' \rangle$. The mapping 
\[
x _{0} \rightarrow f _{0} = J x _{0}
\]
of $X$ into $(X_s')_s'$ satisfies the conditions 
\[
J \left(x _{1} + x _{2}\right) = J x _{1} + J x _{2}, J (\alpha x) = \alpha J (x).
\]
\end{corollaryu}

\begin{theorem}\label{thm:IV.8.Theorem.2} If $X$ is a normed linear space, then the mapping $J$ is isometric, i.e., $\|Jx\| = \|x\|$. 
\end{theorem}

\begin{proof}
 We have $|f_0(x')| = |\langle x_0, x' \rangle| \leqq \|x_0\| \cdot \|x'\|$ so that $\|f_0\| \leqq \|x_0\|$. On the other hand, if $x_0 \neq 0$, then there exists, by Corollary~\hyperref[cor:IV.6.Corollary.2]{$2$} of Theorem~\ref{thm:IV.6.Theorem.1} in Chapter~\ref{chap:IV}, \ref{sec:IV.6}, an element $x_0' \in X'$ such that $\langle x_0, x_0' \rangle = \|x_0\|$ and $\|x_0'\| = 1$. Hence $f_0(x_0') = \langle x_0, x_0' \rangle = \|x_0\|$ so that $\|f_0\| \geqq \|x_0\|$. We have thus proved $\|Jx\| = \|x\|$. 
\end{proof}

\begin{remarku} As the strong dual space of $X_{s}^{\prime}$, the space $(X_{s}^{\prime})_{s}^{\prime}$ is a $B$-space. Hence a normed linear space $X$ may be considered as a linear subspace of the $B$-space $(X_{s}^{\prime})_{s}^{\prime}$ by the embedding $x \to Jx$. Therefore the strong closure of $JX$ in the $B$-space $(X_{s}^{\prime})_{s}^{\prime}$ gives a concrete construction of the completion of $X$. 
\end{remarku}

\begin{definition}\label{def:IV.8.Definition.4} $A$ normed linear space $X$ is said to be reflexive if $X$ may be identified with its second dual or the bidual $(X_{s}^{\prime})_{s}^{\prime}$ by the correspondence $x \leftrightarrow Jx$ above. We know already (see Chapter~\ref{chap:III}, \ref{sec:III.6}) that a Hilbert space is reflexive. As remarked above, $(X_{s}^{\prime})_{s}^{\prime}$ is a $B$-space and so any reflexive normed linear space must be a $B$-space. 
\end{definition}

\begin{theorem}\label{thm:IV.8.Theorem.3} Let $X$ be a $B$-space and $x_0''$ any bounded linear functional on $X_s'$. Then, for any $\varepsilon > 0$ and any finite system of elements $f_1, f_2, \ldots, f_n$ of $X_s'$, there exists an $x_0 \in X$ such that 
\[
\left\| x _{0} \right\| \leqq \left\| x _{0} ^{\prime \prime} \right\| + \varepsilon \quad \text{ and } \quad f _{j} (x _{0}) = x _{0} ^{\prime \prime} (f _{j}) (j = 1, 2, \dots , n).
\]
\end{theorem}

\begin{proof}
 We apply Helly's theorem $5$ in Chapter~\ref{chap:IV}, \ref{sec:IV.6}. For any system of numbers $\beta_{1}, \beta_{2}, \ldots, \beta_{n}$, we have 
\begin{equation}\left| \sum\limits_{j=1} ^{n} \beta_{j} \alpha_{j} \right| = \left| \sum\limits_{j=1} ^{n} \beta_{j} x _{0} ^{\prime \prime} (f _{j}) \right| = \left| x _{0} ^{\prime \prime} \left(\sum\limits_{j=1} ^{n} \beta_{j} f _{j}\right) \right| \leqq \gamma \cdot \left\| \sum\limits_{j=1} ^{n} \beta_{j} f _{j} \right\|, \text{ wehre } \gamma = \| x _{0} ^{\prime \prime} \|, \alpha_{j} = x _{0} ^{\prime \prime} (f _{j}), \notag\end{equation}
and hence, again by Helly's theorem, we obtain an $x_0 \in X$ with the estimate $\| x_0 \| \leqq \gamma + \varepsilon = \| x_0'' \| + \varepsilon$ and $x_0(f_j) = \alpha_j (j = 1, 2, \ldots, n)$. 
\end{proof}

\begin{corollaryu} The unit sphere $S = \{x \in X; \| x \| \leqq 1\}$ of a $B$-space $X$ is dense in the unit sphere of $(X_s')_s'$ in the weak* topology of $(X_s')'$. 
\end{corollaryu}


\yossection{9}{Examples of Dual Spaces}
\label{sec:IV.9}

\begin{example}\label{ex:IV.9.Example.1} $(c_0)' = (l^1)$. 
\end{example}

To any $f \in (c_0)'$, there corresponds a uniquely determined $y_f = \{\eta_n\} \in (\mathcal{I}^1)$ such that, for all $x = \{\xi_n\} \in (c_0)$, 
\begin{equation}\label{eq:IV.9.1}
\langle x, f \rangle = \sum\limits_{n=1} ^{\infty} \xi_{n} \eta_{n} \text{ and } \| f \| = \| y _{f} \|.
\end{equation}

And conversely, any $y = \{\eta_n\} \in (l^1)$ defines an $f_y \in (c_0)'$ such that, for any $x = \{\xi_n\} \in (c_0)$, 
\begin{equation}\tag{1'}\label{eq:IV.9.1'}
\langle x, f _{y} \rangle = \sum\limits_{n=1} ^{\infty} \xi_{n} \eta_{n} \text{ and } \| f _{y} \| = \| y \|.
\end{equation}

\begin{proof}
 Let us define the unit vector $e_k$ by 
\[
e _{k} = \overbrace {(0 , 0 , \dots , 0} ^{k-1}, 1, 0, 0, \dots) (k = 1, 2, \dots). \]

For any $x = \{\xi_n\} \in (c_0)$ and $f \in (c_0)'$, we have, by $s\text{- }\lim\limits_{k \to \infty} \sum\limits_{n=1}^{k} \xi_n e_n = x$. 
\[
\langle x, f \rangle = \lim\limits _{k \rightarrow \infty} \left\langle \sum\limits_{n=1} ^{k} \xi_{n} e _{n}, f \right\rangle = \lim\limits _{k \rightarrow \infty} \sum\limits_{n=1} ^{k} \xi_{n} \eta_{n}, \eta_{n} = f (e _{n}).
\]

Let $\eta_n = \varepsilon_n |\eta_n|$ for $\eta_n \neq 0$, and $\varepsilon_n = \infty$ for $\eta_n = 0$. Take $x^{(n_0)} = \{\xi_n\} \in (c_0)$ in such a way that $\xi_n = \varepsilon_n^{-1}$ for $n \leqq n_0$, and $\xi_n = 0$ for $n > n_0$. Then $\|x^{(n_0)}\| \leqq 1$ and so $\|f\| = \sup\limits_{\|x\| \leqq 1} |\langle x, f \rangle| \geqq |\langle x^{(n_0)}, f \rangle| = \sum\limits_{n=1}^{n_0} |\eta_n|$. Thus, by letting $n_0 \to \infty$, we see that $y_f = \{\eta_n\} \in (l^1)$ and $\|y_f\| = \sum\limits_{n=1}^{\infty} |\eta_n| \leqq \|f\|$. 

If, conversely, $y = \{\eta_n\} \in (l^1)$, then $\left|\sum\limits_{n=1}^{\infty} \xi_n \eta_n\right| \leqq \|x\| \cdot \|y\|$ for all $x = \{\xi_n\} \in (c_0)$, and so $y$ defines an $f_y \in (c_0)'$ and $\|f_y\| \leqq \|y\|$. 
\end{proof}

\begin{example}\label{ex:IV.9.Example.2} $(c)' = (l^1)$. 
\end{example}

For any $x = \{\xi_n\} \in (c)$, we have the representation 
$x = \xi_0e_0 + s\text{- }\lim\limits_{k\to \infty}\sum\limits_{n=1}^{k}(\xi_n - \xi_0)e_n,$ where $\xi_0 = \lim\limits_{n\to \infty}\xi_n,e_0 = (1,1,1,\ldots)$. Thus, for any $f\in (c)^{\prime}$, we have 
\begin{equation}\label{eq:IV.9.2}
\langle x, f \rangle = \xi_{0} \left\langle e _{0}, f \right\rangle + \lim\limits _{k \rightarrow \infty} \left\langle \sum\limits_{n=1} ^{k} \left(\xi_{n} - \xi_{0}\right) e _{n}, f \right\rangle = \xi_{0} \eta_{0} ^{\prime} + \sum\limits_{n=1} ^{\infty} \left(\xi_{n} - \xi_{0}\right) \eta_{n},
\end{equation}
where $\eta_0' = \langle e_0, f \rangle$ and $\eta_n = \langle e_n, f \rangle$ ( $n = 1, 2, \ldots$). As above, we may take $x^{(n_0)} = \{\xi_n\} \in (c_0) \subseteq (c)$ which satisfies 
\[
\left\| x ^ {\left(n _{0}\right)} \right\| \leqq 1, \xi_{0} = \lim\limits _{n \rightarrow \infty} \xi_{n} = 0 \text{ and } \langle x ^ {\left(n _{0}\right)}, f \rangle = \sum\limits_{n=1} ^ {n _{0}} | \eta_{n} |.
\]

Hence, by $|\langle x^{(n_0)}, f \rangle| \leqq \|x^{(n_0)}\| \|f\|$, we see that $\{\eta_n\}^\infty \in (l^1)$. We set $\eta_0' - \sum\limits_{n=1}^\infty \eta_n = \eta_0$. Then, by \eqref{eq:IV.9.2}, we have 
$\langle x, f \rangle = \xi_0 \eta_0 + \sum\limits_{n=1}^{\infty} \xi_n \eta_n$, where $x = \{\xi_n\} \in (c)$ and $\xi_0 = \lim\limits_{n \to \infty} \xi_n$. \eqref{eq:IV.9.2} Let $\eta_n = \varepsilon_n | \eta_n|$ for $\eta_n \neq 0$, and $\varepsilon_n = \infty$ for $\eta_n = 0$ ( $n = 0, 1, 2, \ldots$). Take $x = \{\xi_n\} \in (c)$ such that 
$\xi_{n} = \varepsilon_{n}^{-1}$ if $n\leqq n_0$, and $\xi_{n} = \varepsilon_{0}^{-1}$ if $n > n_0$ 

Then $\|x\| \leqq 1$, $\xi_0 = \lim\limits_{n \to \infty} \xi_n = \varepsilon_0^{-1}$, and $\langle x, f \rangle = |\eta_0| + \sum\limits_{n=1}^{n_0} |\eta_n| + \varepsilon_0^{-1} \sum\limits_{n=n_0+1}^{\infty} \eta_n$. Hence, we must have $|\eta_0| + \sum\limits_{n=1}^{\infty} |\eta_n| \leqq \|f\|$. 

If, conversely, $y = \{\eta_n\}_0^\infty$ is such that $\|y\| = |\eta_0| + \sum\limits_{n=1}^{\infty} |\eta_n| < \infty$, then 
\[
\eta_{0} \cdot \lim\limits _{n \rightarrow \infty} \xi_{n} + \sum\limits_{n=1} ^{\infty} \xi_{n} \eta_{n}, \text{ where } x = \{\xi_{n} \} _{1} ^{\infty} \in (c),
\]
defines an $f_y \in (c)'$ such that $\|f_y\| \leqq |\eta_0| + \sum\limits_{n=1}^{\infty} |\eta_n|$. 

Therefore, we have proved that $(c)' = (l^1)$ as explained above. 

\begin{example}\label{ex:IV.9.Example.3} $L^p(S, \mathfrak{B}, m)' = L^q(S, \mathfrak{B}, m)$ ( $1 \leqq p < \infty$ and $p^{-1} + q^{-1} = 1$). To any $f \in L^p(S)'$, there corresponds a $y_f \in L^q(S)$ such that 
\begin{equation}\label{eq:IV.9.3}
\langle x, f \rangle = \int_{S} x (s) y _{f} (s) m (\dd s) \text{ for all } x \in L ^{p} (S) \text{ and } \| f \| = \| y _{f} \|,
\end{equation}
and conversely, any $y \in L^{q}(S)$ defines an $f_{y} \in L^{p}(S)'$ such that 
\begin{equation}\tag{3'}\label{eq:IV.9.3'}
\langle x, f _{y} \rangle = \int_{S} x (s) y (s) m (\dd s) \text{ for all } x \in L ^{p} (S) \text{ and } \| f _{y} \| = \| y \|.
\end{equation}
\end{example}

\begin{proof}
 Let $S = \bigcup\limits_{j=1}^{\infty} B_j$ with $0 < m(B_j) < \infty$ and set $B^{(n)} = \bigcup\limits_{j=1}^{n} B_j$. For a fixed $n$, the defining function $C_B(s)$ of the set $B \subseteq B^{(n)}$ is $\in L^p(S)$. Thus the set function $\psi(B) = \langle C_B, f \rangle$ is $\sigma$-additive and $m$-absolutely continuous in $B \subseteq B^{(n)}$. By the differentiation theorem of Lebesgue-Nikodym (see Chapter~\ref{chap:III}, \ref{sec:III.8}), there exists a $y_n(s) \in L^1(B^{(n)}, \mathfrak{B}^{(n)}, m)$ such that 
\[
\psi (B) = \int_{B} y _{n} (s) m (\dd s) \text{ whenever } B \subseteq B ^{(n)},
\]
the family $\mathfrak{B}^{(n)}$ of sets being defined by $\mathfrak{B}^{(n)} = \{B \cap B^{(n)}; B \in \mathfrak{B}\}$. Therefore, by setting $y(s) = y_n(s)$ for $s \in B^{(n)}$, we have 
\[
\langle C _{B}, f \rangle = \int_{B} y (s) m (\dd s) \text{ for } B \in \mathfrak {B} ^{(n)} (n = 1, 2, \dots). \]

Hence, for any finitely-valued function $x$ with support in some $B^{(n)}$, 
\begin{equation}\label{eq:IV.9.4}
\langle x, f \rangle = \int_{S} x (s) y (s) m (\dd s).
\end{equation}

Let $x \in L^{p}(S)$ and put 
\[
\begin{aligned} x _{n} (s) & = x (s) \text{ if } | x (s) | \leqq n \text{ and } s \in B ^{(n)}, \\ & = 0 \text{ otherwise }. \end{aligned}
\]

We decompose the set $\{z; |z| \leqq n\}$ of the complex plane into a finite number of disjoint Baire sets $M_{n,k,t}$ ( $t = 1, 2, \ldots, d_{k,n}$) of diameters $\leqq 1/k$. Set, for the $x_n(s) \in L^\infty(S, \mathfrak{B}, m)$, 
$x_{n,k}(s) =$ a constant $z$, such that $z\in$ (the closure $M_{n,k,t}^{a}$) and $|z| = \inf\limits_{w\in Mn,k,t}|w|$ whenever $x_{n}(s)\in M_{n,k,t}$. 

Then $|x_{n,k}(s)| \leqq |x_n(s)|$ and $\lim\limits_{k \to \infty} x_{n,k}(s) = x_n(s)$ and so, by the Lebesgue-Fatou Lemma, $\slim\limits_{k \to \infty} x_{n,k} = x_n (n = 1, 2, \ldots)$. Thus, again by the Lebesgue-Fatou Lemma, 
\begin{equation}\label{eq:IV.9.5}
\begin{aligned} \langle x _{n}, f \rangle & = \lim\limits _{k \to \infty} \langle x _{n,k}, f \rangle = \lim\limits _{k \to \infty} \int_{S} x _{n,k} (s) y (s) m (\dd s) \\ & = \int_{S} \lim\limits _{k \to \infty} x _{n,k} (s) \cdot y (s) m (\dd s) = \int_{S} x _{n} (s) y (s) m (\dd s). \end{aligned}
\end{equation}

Since \slim $x_{n} = x$, we see that $\langle x, f \rangle = \lim\limits_{n \to \infty} \langle x_{n}, f \rangle$. We put, for any complex number $z$, $a(z) = e^{-i\theta}$ if $z = r e^{i\theta}$ and $a(0) = 0$. Then $\|x\| \geqq \|(|x_{n}| \cdot a(y))\|$ and so 
\[
\left\| f \right\| \left\| x \right\| \geqq \langle \left| x _{n} \right| a (y), f \rangle = \int_{S} \left| x _{n} (s) \right| \cdot | y (s) | m (\dd s).
\]

Thus, by the Lebesgue-Fatou Lemma, $\|f\| \cdot \|x\| \geqq \int_{S} |x(s)\|y(s)| m(\dd s)$ and so the function $x(s)y(s)$ belongs to $\in L^{1}(S)$. Therefore, letting 
$n\to \infty$ in \eqref{eq:IV.9.5}, we obtain 
\[
\langle x, f \rangle = \int_{S} x (s) y (s) m (\dd s) \text{ whenever } x \in L ^{p} (S).
\]

We shall show that $y \in L^q(S)$. To this purpose, set 
\[
\begin{aligned} y _{n} (s) & = y (s) \text{ if } | y (s) | \leqq n \text{ and } s \in B ^{(n)}, \\ & = 0 \text{ otherwise }. \end{aligned}
\]

Then $y_{n} \in L^{q}(S)$ and, as proved above, 
\[
\begin{aligned}
&\left\| f \right\| \cdot \left\| x \right\|  \geqq \langle | x | \cdot a (y), f \rangle = \int_{S} | x (s) | | y (s) | m (\dd s) \geqq \int_{S} | x (s) | | y _{n} (s) | m (\dd s).
\end{aligned}
\]

If we take $x(s) = |y_n(s)|^{q/p}$ and apply Hölder's equality, we obtain 
\[
\int_{S} | x (s) | | y _{n} (s) | m (\dd s) = \left(\int_{S} | x (s) | ^{p} m (\dd s)\right) ^{1/p} \left(\int_{S} | y _{n} (s) | ^{q} m (\dd s)\right) ^{1/q}.
\]

Hence $\|f\| \geqq \|y_n\| = \left( \int_S |y_n(s)|^q m(\dd s) \right)^{1/q}$, with the understanding that, when $p = 1$ we have $\|f\| \geqq \|y_n\| = \text{ essential sup }_{s \in S} |y_n(s)|$. 

Therefore, by letting $n \to \infty$ and applying the Lebesgue-Fatou Lemma, we see that $y \in L^{q}(S)$ and $\| f \| \geqq \| y \|$. On the other hand, any $y \in L^{q}(S)$ defines an $f \in L^{p}(S)'$ by $\langle x, f \rangle = \int_{S} x(s) y(s) m(\dd s)$ as may be seen by Hölder's inequality, and this inequality shows that $\| f \| \leqq \| y \|$. 
\end{proof}

\begin{remarku} We have incidentally proved that 
\[
(l ^{p}) ^{\prime} = (l ^{q}) (1 \leqq p < \infty \text{ and } p ^{-1} + q ^{-1} = 1).
\]
\end{remarku}

\begin{example}\label{ex:IV.9.Example.4} Let the measure space $(S, \mathfrak{B}, m)$ with $m(S) < \infty$ have the property that, for any $B \in \mathfrak{B}$ with $0 < m(B) = \delta < \infty$ and positive integer $n$, there exists a subset $B_n$ of $B$ such that $\delta(n + 1)^{-1} \leqq m(B_n) \leqq \delta n^{-1}$. Then no other continuous linear functionals $\in M(S, \mathfrak{B}, m)$ ' than the zero functional can exist. 
\end{example}

\begin{proof}
 Any $x \in L^{1}(S, \mathfrak{B}, m)$ belongs to $M(S, \mathfrak{B}, m)$ and the topology of $L^{1}(S, \mathfrak{B}, m)$ is stronger than that of $M(S, \mathfrak{B}, m)$. Thus any $f \in M(S, \mathfrak{B}, m)'$, when restricted to the functions of $L^{1}(S, \mathfrak{B}, m)$, defines a continuous linear functional $f_{0} \in L^{1}(S, \mathfrak{B}, m)'$. Thus there exists a $y \in L^{\infty}(S, \mathfrak{B}, m)$ such that 
\[
\langle x, f \rangle = \langle x, f _{0} \rangle = \int_{S} x (s) y (s) m (\dd s) \text{ whenever } x \in L ^{1} (S, \mathfrak {B}, m). \]

Since $L^1(S, \mathfrak{B}, m)$ is dense in $M(S, \mathfrak{B}, m)$ in the topology of $M(S, \mathfrak{B}, m)$, the condition $f \neq 0$ implies that $f_0 \neq 0$. Thus there exists an $\varepsilon > 0$ such that $B = \{s; |y(s)| \geqq \varepsilon\}$ has its measure $m(B) = \delta > 0$. Let $B_n \subseteq B$ be as in the hypothesis, and let $y(s) = r e^{i\theta}$ for $s \in B$. Set $x_{n}(s) = e^{-i\theta}$ for $s \in B_{n}$ and $x_{n}(s) = 0$ otherwise. Then $z_{n}(s) = nx_{n}(s)$ converges to $0$ asymptotically, that is, \slim $z_{n} = 0$ in $M(S, \mathfrak{B}, m)$. But 
\[
\lim\limits _{n \rightarrow \infty} \langle z _{n}, f \rangle = \lim\limits _{n \rightarrow \infty} \langle z _{n}, f _{0} \rangle = \lim\limits _{n \rightarrow \infty} \int_{S} z _{n} (s) y (s) m (\dd s) \geqq \delta \varepsilon > 0,
\]
contrary to the continuity of the functional $f$. 
\end{proof}

\begin{example}\label{ex:IV.9.Example.5} $L^{\infty}(S, \mathfrak{B}, m)'$. 
\end{example}

Let an $f \in L^{\infty}(S, \mathfrak{B}, m)'$ be given and set, for any $B \in \mathfrak{B}$, $f(C_B) = \psi(B)$ where $C_B(s)$ is the defining function of the set $B$. We have then: 
\begin{equation}\label{eq:IV.9.6}
B _{1} \cap B _{2} = \emptyset \text{ implies } \psi (B _{1} + B _{2}) = \psi (B _{1}) + \psi (B _{2}),
\end{equation}
that is $\psi$ is finitely additive, 
the real part $\psi_{1}(B)$ and the imaginary part $\psi_{2}(B)$ of 
$\psi (B)$ are of bounded total variation, that is, $\sup\limits_B |\psi_i(B)|$ 
\begin{equation}\label{eq:IV.9.7}
< \infty (i = 1, 2),
\end{equation}
$\psi$ is $m$-absolutely continuous, that is $m(B) = 0$ implies 
\begin{equation}\label{eq:IV.9.8}
\psi (B) = 0.
\end{equation}

The condition \eqref{eq:IV.9.6} is a consequence of the linearity of $f$, and \eqref{eq:IV.9.7} and \eqref{eq:IV.9.8} are clear from $|\psi(B)| \leqq \|f\| \cdot \|C_B\|$. 

For any $x \in L^{\infty}(S, \mathfrak{B}, m)$, we consider a partition of the sphere $\{z; |z| \leqq \|x\|\}$ of the complex plane into a finite system of disjoint Baire sets $A_{1}, A_{2}, \ldots, A_{n}$ of diameters $\leqq \varepsilon$. If we set $B_{i} = \{s \in S; x(s) \in A_{i}\}$, then, no matter what point $\alpha_{i}$ we choose from $A_{i} (i = 1, 2, \ldots, n)$, we 
\[
\left\| x - \sum\limits_{i=1} ^{n} \alpha_{i} C _ {B _{i}} \right\| \leqq \varepsilon ,
\]
and so 
\[
\left| f (x) - \sum\limits_{i=1} ^{n} \alpha_{i} \psi (B _{i}) \right| \leqq \| f \| \cdot \varepsilon .
\]

Thus, if we let $n \to \infty$ in such a way that $\varepsilon \downarrow 0$, we obtain 
\begin{equation}\label{eq:IV.9.9}
f (x) = \lim \sum_{i=1}^{n} \alpha_{i} \psi (B _{i}),
\end{equation}
independently of the manner of partition $\{z; |z| \leqq \|x\|\} = \sum\limits_{i=1}^{n} A_i$ and choice of points $\alpha$'s. The limit on the right of \eqref{eq:IV.9.9} is called Radon's integral of $x(s)$ with respect to the finitely additive measure $\psi$. Thus 
\begin{equation}\label{eq:IV.9.10}
f(x) = \int_{S} x(s) \psi(\dd s) \text{ (Radon's integral) whenever } x \in L^{\infty}(S, \mathfrak{B}, m),
\end{equation}
and so 
\begin{equation}\label{eq:IV.9.11}
\left\| f \right\| = \sup\limits _ {\text{ ess. } \sup | x (s) | \leqq 1} \left| \int_{S} x (s) \psi (\dd s) \right|.
\end{equation}

Conversely, it is easy to see that any $\psi$ satisfying \eqref{eq:IV.9.6}, \eqref{eq:IV.9.7} and \eqref{eq:IV.9.8} defines an $f \in L^{\infty}(S, \mathfrak{B}, m)'$ through \eqref{eq:IV.9.10} and that \eqref{eq:IV.9.11} is true. 

Therefore, we have proved that $L^{\infty}(S,\mathfrak{B},m)'$ is the space of all set functions $\psi$ satisfying \eqref{eq:IV.9.6}, \eqref{eq:IV.9.7} and \eqref{eq:IV.9.8} and normed by the right hand side of \eqref{eq:IV.9.11}, the so-called total variation of $\psi$. 

\begin{remarku} We have so far proved that $L^p(S, \mathfrak{B}, m)$ is reflexive when $1 < p < \infty$. However, the space $L^1(S, \mathfrak{B}, m)$ is, in general, not reflexive. 
\end{remarku}

\begin{example}\label{ex:IV.9.Example.6} $C(S)'$.
\end{example}

Let $S$ be a compact topological space. Then the dual space $C(S)'$ of the space $C(S)$ of complex-valued continuous functions on $S$ is given as follows. To any $f \in C(S)'$, there corresponds a uniquely determined complex Baire measure $\mu$ on $S$ such that 
\begin{equation}\label{eq:IV.9.12}
f (x) = \int_{S} x (s) \mu (\dd s) \text{ whenever } x \in C (S),
\end{equation}
and hence 
\begin{equation}\label{eq:IV.9.13}
\left\| f \right\| = \sup\limits _ {\sup\limits _{s} | x (s) | \leqq 1} \left| \int_{S} x (s) \mu (\dd s) \right| = \text{ the total variation of } \mu \text{ on } S.
\end{equation}

Conversely, any Baire measure $\mu$ on $S$ such that the right side of \eqref{eq:IV.9.13} is finite, defines a continuous linear functional $f \in C(S)'$ through \eqref{eq:IV.9.12} and we have \eqref{eq:IV.9.13}. Moreover, if we are concerned with a real functional $f$ on a real $B$-space $C(S)$, then the corresponding measure $\mu$ is real-valued; if, moreover $f$ is positive, in the sense that $f(x) \geqq 0$ for non-negative functions $x(s)$, then the corresponding measure $\mu$ is positive, i.e., $\mu(B) \geqq 0$ for every $B \in \mathfrak{B}$. 

\begin{remarku} The result stated above is known as the F. Riesz-A. Markov-S. \textsc{Kakutani} theorem, and is one of the fundamental theorems in topological measures. For the proof, the reader is referred to standard text books on measure theory, e.g., P. R. \textsc{Halmos} \cite{halmos92} and N. \textsc{Dunford}-J. \textsc{Schwartz} \cite{dunford44}. 
\end{remarku}

\yossub{References for Chapter IV}

For the Hahn-Banach theorems and related topics, see \textsc{Banach} \cite{banach9}, \textsc{Bourbaki} \cite{bourbaki27} and \textsc{Köthe} \cite{konishi164}. It was \textsc{Mazur} \cite{mazur192} who noticed the importance of convex sets in normed linear spaces. The proof of Helly's theorem given in this book is due to Y. \textsc{Mimura} (unpublished). 
